\chapter{Primene}
Dokazi nultog znanja obuhvataju \v sirok spektar slu\v cajeva kori\v s\' cenja i uveliko se koriste u svakodnevnom životu, ali postoji još prostora za rast i razvoj primena ove tehnologije u budu\' cnosti. U nastavku su navedeni neki od naj\v ce\v s\' ci slu\v cajevi kori\v s\' cenja dokaza nultog znanja.

\section{Anonimna pla\' canja}
 Najrasprostranjeniji oblik plaćanja danas je plaćanje kreditnom. Plaćanja kreditnom karticom su obi\v cno vidljiva banci, vladi i nekim drugim zainteresovanim stranama. Što mo\v ze biti korisno pri identifikaciji ilegalnih aktivnosti, ali sa druge strane ugro\v zava privatnost obi\v cnih ljudi. Kako bi se obezbedile u potpunosti privatne transakcije napravljene su kriptovalute. Me\dj utim ve\' cina transakcija kriptovalutama je vidljiva na javnim blokčejnovima, gde se ponovo mogu vršiti restrikcije korisničke privatnosti. Korisni\v cki identiteti su \v cesto pseudonimni ili su svojevoljno povezani sa stvarnim identitetima (na primer \textenglish{GitHub} profilima) ali se takođe analizom podataka mogu povezati sa stvarnim identitetima.\\

Da bi se obezbedila maksimalna privatnost, napravljeni su specijalni \textenglish{privacy coins} koji predstavljaju tip kriptovaluta koji je  dizajniran sa akcentom na privatnost i anonimnost korisnika. On \v stiti detalje transakcije, uklju\v cuju\' ci adrese primaoca i po\v siljaoca, vrstu i koli\v cinu sredstava kao i vreme transakcije. Najpopularnije kriptovalute ovog tipa su \textenglish{Monero}\cite{monero}, \textenglish{Zcash, Horizen}\cite{horizen}, \textenglish{Aleo}\cite{aleo} i \textenglish{Dash}\cite{Duffield2017DashAP}, i upravo one u svojim transakcijama koriste dokaze nultog znanja. Postoje i razni protokoli zasnovani na dokazima nultog znanja koji doprinose anonimnim plaćanjima. Najpoznatiji je \textenglish{Aztec} protokol\cite{aztec} koji omogućava privatne transakcije na \textenglish{Ethereum} mreži.

\section{Autentifikacija}
Dokazi nultog znanja pronalaze primenu i u autentifikacionim protokolima. Najvi\v se se koriste u autentifikacionim sistemima zasnovanim na lozinkama. Oni dozvoljavaju korisniku da doka\v ze znanje tajne \v sifre bez potrebe da stvarno otkrije \v sifru.\\

\textenglish{Sigma} protokoli su klasični primer interaktivnih protokola nultog znanja koji se koriste za autentifikaciju\cite{inbook}. Izvršavaju se iz 3 koraka: izazov, odgovor i verifikacija. U prvom koraku se zadaje izazov na koji se u drugom odgovara uz pomoć korišćenja tajne informacije ali bez njenog otkrivanja. U poslednjem koraku se odgovor verifikuje.\\

\textenglish{ZPass} predstavlja još jedan primer protokola sa nultim znanjem za verifikovanje identiteta. On pojednostavljuje verifikaciju državljanstva globalnim preduzećima obezbeđujući usklađenost sa internacionalnim zakonima. Pored toga nudi i verifikaciju godina korisnika za internet prodavnice alkohola i industrije video igara\cite{zpass}. Pored ovog protokola, poznati protokoli za svrhe verifikovanja identiteta su \textenglish{Semaphore}\cite{semaforrad, semwrl} i \textenglish{Galxe protocol}\cite{galxe}. 



\section{Za\v stita identiteta}
Kada je potrebno dokazati da korisnik poseduje odre\dj ene informacije bez otkrivanja samih informacija mogu se koristiti dokazi nultog znanja. Na primer ako korisnik \v zeli da doka\v ze da je punoletan ili da ima odre\dj eno dr\v zavljanstvo. \\

Pored toga primena dokaza nultog znanja se pronalazi i u zdravstvu. U zdravstvenim aplikacijama radi se sa osetljivim informacijama o pacijentima, pa ako pacijent \v zeli da doka\v ze da ispunjava uslove za odre\dj enu medicinsku proceduru, mo\v ze to posti\' ci bez otkrivanja samih podataka o sebi. Primer ovakvog sistema je \textenglish{Health-zkIDM}. To je decentralizovani sistem autentifikacije identiteta koji je zasnovan na dokazima nultog znanja i tehnologiji blokčejna. On omogućava pacijentima da se identifikuju i verifikuju svoje identitete transparentno i bezbedno\cite{s22207716}. 

\section{Pametni ugovori u blokčejn tehnologiji}
Pametni ugovori u blokčejnu su samostalni programski kodovi u vidu digitalnih ugovora koji se automatski izvršavaju ukoliko su ispunjeni unapred zadati uslovi. Oni primenjuju dokaze nultog znanja kako bi se omogu\' cila provera odre\dj enih uslova bez otkrivanja detalja ugovora ili podataka koji se koriste u transakcijama. Tako je omogu\' cena sigurnost i privatnost u okviru decentralizovanih platformi. Na primer, u decentralizovanim finansijama korisnici mogu dokazati da ispunjavaju uslove za podizanje određenog kredita bez otkrivanja svojih finansijskih podataka\cite{inMaobook}.

\section{Potencijalne primene}
Dokazi nultog znanja se u budu\' cnosti mogu primenjivati na razli\v cite oblasti poput obrazovanja, gde bi studentima bila zagarantovana ve\' ca privatnost. Autenti\v cnost diploma bi se mogla potvrditi bez otkrivanja li\v cnih podataka. Kada bi na primer student \v zeleo da aktivira besplatno profesionalnu verziju {\it \textenglish{GitHub}} naloga za koju kao student ima pravo,  ne bi morao da \v salje fotografiju identifikacionog dokumenta sa obele\v zjem fakulteta kojim bi izlo\v zio dodatne informacije o sebi. Ve\' c bi  mogao pomo\' cu nekog protokola nultog znanja da izgeneri\v se dokaz koji bi \textenglish{GitHub} mogao da verifikuje. Time \textenglish{GitHub} ne bi znao na kom fakultetu studira student ili bilo kakve propratne informacije o njemu koje bi dobio slikom dokumenta, ali bi bio ube\dj en da je on student. Popularizacijom dokaza nultog znanja, njihova primena bi se mogla proširiti na razne sfere svakodnevnog života. Dokazi nultog znanja bi se mogli koristiti za novi sistem elektronskog glasanja u kome bi pru\v zili sigurnu verifikaciju identiteta bira\v ca i korektnost glasanja bez otkrivanja pojedina\v cnih glasova.